\documentclass[../main.tex]{subfiles}

\begin{document}

Tomamos prestado el lenguaje \texttt{IMP} de~\cite{winskel1993formal}.
Es un pequeño lenguaje \enquote{imperativo}, en el sentido de que tiene variables, que 
conservan el \emph{estado} del programa, y \emph{comandos} que modifican
este estado.

Presentamos un ejemplo de programa en \texttt{IMP} que calcula el factorial de un número \texttt{n}:

\begin{minted}{c}
n = 5; i = 2; res = 1;

while i <= n {
    res = res * i;
    i = i + 1;
}
\end{minted}

Utilizaremos Lean como nuestro \enquote{metalenguaje} para definir la sintaxis y la semántica
de \texttt{IMP}.

\section{Sintaxis}

Para definir la sintaxis de \texttt{IMP} utilizaremos los tipos inductivos de Lean,
de este modo obtendremos primero el \enquote{árbol de sintaxis abstracta} del lenguaje.
Distinguimos entre expresiones aritméticas, expresiones booleanas y comandos.
La distinción intuitiva entre expresiones y comandos es que las primeras
\enquote{devuelven} un valor, mientras que los segundos \enquote{modifican} el estado.
Veremos que en realidad, esto no es del todo cierto.

Comenzamos por la expresiones aritméticas, solo dependen de \texttt{Int} y \texttt{String}.

\begin{minted}{lean4}
inductive Aexp where
  | val : Int → Aexp
  | var : String → Aexp
  -- arithmetic
  | add : Aexp → Aexp → Aexp
  | sub : Aexp → Aexp → Aexp
  | mul : Aexp → Aexp → Aexp
\end{minted}

Podemos interpretar la anterior definición como sigue, el tipo \texttt{Aexp} tiene
cinco constructores, a los que Intuitivamente asociamos los siguientes significados:
\begin{itemize}
  \item [\texttt{val}:] construye la expresión formada por una única constante entera.
  \item [\texttt{var}:] construye la expresión formada por una única variable.
  \item [\texttt{add}:] construye la expresión formada por la suma de dos expresiones aritméticas.
  \item [\texttt{sub}:] construye la expresión formada por la resta de dos expresiones aritméticas.
  \item [\texttt{mul}:] construye la expresión formada por la multiplicación de dos expresiones aritméticas.
\end{itemize}

En realidad el \enquote{constructor de suma} \texttt{add} solo es \enquote{de suma} en
el nombre, ya que hasta que no definamos una semántica para las expresiones
aritméticas lo único que expresa este es que se puede formar una
nueva expresión a partir de dos expresiones ya existentes.

Lo mismo podemos decir del constructor \texttt{var}, que nos permitirá
\enquote{acceder} al contenido de una variable. En realidad, sintacticamente,
solo nos permite \enquote{construir} una expresión aritmética a partir de un \texttt{String}.

Por ejemplo podemos representar la expresión aritmética \(x + 3\) como 
\begin{minted}{lean4}
#check Aexp.add (Aexp.var "x") (Aexp.val 3)
\end{minted}

Seguimos con las expresiones booleanas.

\begin{minted}{lean4}
inductive Bexp where
  -- constants
  | tt
  | ff
  -- boolean
  | not : Bexp → Bexp
  | and : Bexp → Bexp → Bexp
  | or  : Bexp → Bexp → Bexp
  -- comparison
  | eq : Aexp → Aexp → Bexp
  | le : Aexp → Aexp → Bexp
\end{minted}

Tenemos los siguientes constructores:
\begin{itemize}
  \item [\texttt{tt}:] representa el valor de verdad \(\top\).
  \item [\texttt{ff}:] representa el valor de verdad \(\bot\).
  \item [\texttt{not}:] construye la expresión booleana formada por la negación de otra expresión booleana.
  \item [\texttt{and}:] construye la expresión booleana formada por la conjunción de dos expresiones booleanas.
  \item [\texttt{or}:] construye la expresión booleana formada por la disyunción de dos expresiones booleanas.
  \item [\texttt{eq}:] construye la expresión booleana formada por la comparación de igualdad entre dos expresiones aritméticas.
  \item [\texttt{le}:] construye la expresión booleana formada por la comparación de menor o igual entre dos expresiones aritméticas.
\end{itemize}

Por ejemplo, podemos representar la expresión booleana \(!(x \le 3)\) como
\begin{minted}{lean4}
#check Bexp.not (Bexp.le (Aexp.var "x") (Aexp.val 3))
\end{minted}

Pasamos al tipo más complicado, los comandos. Esto se debe a que todo programa 
de \text{IMP} tendrá tipo \texttt{Com}, en efecto, todos los programas son una 
concatenación de comandos.

\begin{minted}{lean4}
inductive Com where
  | skip
  | cat   : Com → Com → Com
  | ass   : String → Aexp → Com
  | cond  : Bexp → Com → Com → Com
  | loop  : Bexp → Com → Com
\end{minted}

Tenemos los constructores:
\begin{itemize}
  \item[\texttt{skip}:] construye el programa vacío.
  \item[\texttt{cat}:] concatena comandos para formar uno nuevo.
  \item[\texttt{ass}:] corresponde a la asignación de expresiones a variables.
  \item[\texttt{cond}:] corresponde al constructo de flujo de control \texttt{if}.
\end{itemize}

\end{document}